% Options for packages loaded elsewhere
% Options for packages loaded elsewhere
\PassOptionsToPackage{unicode}{hyperref}
\PassOptionsToPackage{hyphens}{url}
%
\documentclass[
  ignorenonframetext,
  aspectratio=169,
]{beamer}
\newif\ifbibliography
\usepackage{pgfpages}
\setbeamertemplate{caption}[numbered]
\setbeamertemplate{caption label separator}{: }
\setbeamercolor{caption name}{fg=normal text.fg}
\beamertemplatenavigationsymbolsempty
% remove section numbering
\setbeamertemplate{part page}{
  \centering
  \begin{beamercolorbox}[sep=16pt,center]{part title}
    \usebeamerfont{part title}\insertpart\par
  \end{beamercolorbox}
}
\setbeamertemplate{section page}{
  \centering
  \begin{beamercolorbox}[sep=12pt,center]{section title}
    \usebeamerfont{section title}\insertsection\par
  \end{beamercolorbox}
}
\setbeamertemplate{subsection page}{
  \centering
  \begin{beamercolorbox}[sep=8pt,center]{subsection title}
    \usebeamerfont{subsection title}\insertsubsection\par
  \end{beamercolorbox}
}
% Prevent slide breaks in the middle of a paragraph
\widowpenalties 1 10000
\raggedbottom
\AtBeginPart{
  \frame{\partpage}
}
\AtBeginSection{
  \ifbibliography
  \else
    \frame{\sectionpage}
  \fi
}
\AtBeginSubsection{
  \frame{\subsectionpage}
}
\usepackage{iftex}
\ifPDFTeX
  \usepackage[T1]{fontenc}
  \usepackage[utf8]{inputenc}
  \usepackage{textcomp} % provide euro and other symbols
\else % if luatex or xetex
  \usepackage{unicode-math} % this also loads fontspec
  \defaultfontfeatures{Scale=MatchLowercase}
  \defaultfontfeatures[\rmfamily]{Ligatures=TeX,Scale=1}
\fi
\usepackage{lmodern}

\usetheme[]{Boadilla}
\usecolortheme[]{default}
\usefonttheme[]{professionalfonts}
\ifPDFTeX\else
  % xetex/luatex font selection
\fi
% Use upquote if available, for straight quotes in verbatim environments
\IfFileExists{upquote.sty}{\usepackage{upquote}}{}
\IfFileExists{microtype.sty}{% use microtype if available
  \usepackage[]{microtype}
  \UseMicrotypeSet[protrusion]{basicmath} % disable protrusion for tt fonts
}{}
\makeatletter
\@ifundefined{KOMAClassName}{% if non-KOMA class
  \IfFileExists{parskip.sty}{%
    \usepackage{parskip}
  }{% else
    \setlength{\parindent}{0pt}
    \setlength{\parskip}{6pt plus 2pt minus 1pt}}
}{% if KOMA class
  \KOMAoptions{parskip=half}}
\makeatother

\usepackage{color}
\usepackage{fancyvrb}
\newcommand{\VerbBar}{|}
\newcommand{\VERB}{\Verb[commandchars=\\\{\}]}
\DefineVerbatimEnvironment{Highlighting}{Verbatim}{commandchars=\\\{\}}
% Add ',fontsize=\small' for more characters per line
\usepackage{framed}
\definecolor{shadecolor}{RGB}{241,243,245}
\newenvironment{Shaded}{\begin{snugshade}}{\end{snugshade}}
\newcommand{\AlertTok}[1]{\textcolor[rgb]{0.68,0.00,0.00}{#1}}
\newcommand{\AnnotationTok}[1]{\textcolor[rgb]{0.37,0.37,0.37}{#1}}
\newcommand{\AttributeTok}[1]{\textcolor[rgb]{0.40,0.45,0.13}{#1}}
\newcommand{\BaseNTok}[1]{\textcolor[rgb]{0.68,0.00,0.00}{#1}}
\newcommand{\BuiltInTok}[1]{\textcolor[rgb]{0.00,0.23,0.31}{#1}}
\newcommand{\CharTok}[1]{\textcolor[rgb]{0.13,0.47,0.30}{#1}}
\newcommand{\CommentTok}[1]{\textcolor[rgb]{0.37,0.37,0.37}{#1}}
\newcommand{\CommentVarTok}[1]{\textcolor[rgb]{0.37,0.37,0.37}{\textit{#1}}}
\newcommand{\ConstantTok}[1]{\textcolor[rgb]{0.56,0.35,0.01}{#1}}
\newcommand{\ControlFlowTok}[1]{\textcolor[rgb]{0.00,0.23,0.31}{\textbf{#1}}}
\newcommand{\DataTypeTok}[1]{\textcolor[rgb]{0.68,0.00,0.00}{#1}}
\newcommand{\DecValTok}[1]{\textcolor[rgb]{0.68,0.00,0.00}{#1}}
\newcommand{\DocumentationTok}[1]{\textcolor[rgb]{0.37,0.37,0.37}{\textit{#1}}}
\newcommand{\ErrorTok}[1]{\textcolor[rgb]{0.68,0.00,0.00}{#1}}
\newcommand{\ExtensionTok}[1]{\textcolor[rgb]{0.00,0.23,0.31}{#1}}
\newcommand{\FloatTok}[1]{\textcolor[rgb]{0.68,0.00,0.00}{#1}}
\newcommand{\FunctionTok}[1]{\textcolor[rgb]{0.28,0.35,0.67}{#1}}
\newcommand{\ImportTok}[1]{\textcolor[rgb]{0.00,0.46,0.62}{#1}}
\newcommand{\InformationTok}[1]{\textcolor[rgb]{0.37,0.37,0.37}{#1}}
\newcommand{\KeywordTok}[1]{\textcolor[rgb]{0.00,0.23,0.31}{\textbf{#1}}}
\newcommand{\NormalTok}[1]{\textcolor[rgb]{0.00,0.23,0.31}{#1}}
\newcommand{\OperatorTok}[1]{\textcolor[rgb]{0.37,0.37,0.37}{#1}}
\newcommand{\OtherTok}[1]{\textcolor[rgb]{0.00,0.23,0.31}{#1}}
\newcommand{\PreprocessorTok}[1]{\textcolor[rgb]{0.68,0.00,0.00}{#1}}
\newcommand{\RegionMarkerTok}[1]{\textcolor[rgb]{0.00,0.23,0.31}{#1}}
\newcommand{\SpecialCharTok}[1]{\textcolor[rgb]{0.37,0.37,0.37}{#1}}
\newcommand{\SpecialStringTok}[1]{\textcolor[rgb]{0.13,0.47,0.30}{#1}}
\newcommand{\StringTok}[1]{\textcolor[rgb]{0.13,0.47,0.30}{#1}}
\newcommand{\VariableTok}[1]{\textcolor[rgb]{0.07,0.07,0.07}{#1}}
\newcommand{\VerbatimStringTok}[1]{\textcolor[rgb]{0.13,0.47,0.30}{#1}}
\newcommand{\WarningTok}[1]{\textcolor[rgb]{0.37,0.37,0.37}{\textit{#1}}}

\usepackage{longtable,booktabs,array}
\newcounter{none} % for unnumbered tables
\usepackage{calc} % for calculating minipage widths
\usepackage{caption}
% Make caption package work with longtable
\makeatletter
\def\fnum@table{\tablename~\thetable}
\makeatother
\usepackage{graphicx}
\makeatletter
\newsavebox\pandoc@box
\newcommand*\pandocbounded[1]{% scales image to fit in text height/width
  \sbox\pandoc@box{#1}%
  \Gscale@div\@tempa{\textheight}{\dimexpr\ht\pandoc@box+\dp\pandoc@box\relax}%
  \Gscale@div\@tempb{\linewidth}{\wd\pandoc@box}%
  \ifdim\@tempb\p@<\@tempa\p@\let\@tempa\@tempb\fi% select the smaller of both
  \ifdim\@tempa\p@<\p@\scalebox{\@tempa}{\usebox\pandoc@box}%
  \else\usebox{\pandoc@box}%
  \fi%
}
% Set default figure placement to htbp
\def\fps@figure{htbp}
\makeatother





\setlength{\emergencystretch}{3em} % prevent overfull lines

\providecommand{\tightlist}{%
  \setlength{\itemsep}{0pt}\setlength{\parskip}{0pt}}



 


\makeatletter
\@ifpackageloaded{tcolorbox}{}{\usepackage[skins,breakable]{tcolorbox}}
\@ifpackageloaded{fontawesome5}{}{\usepackage{fontawesome5}}
\definecolor{quarto-callout-color}{HTML}{909090}
\definecolor{quarto-callout-note-color}{HTML}{0758E5}
\definecolor{quarto-callout-important-color}{HTML}{CC1914}
\definecolor{quarto-callout-warning-color}{HTML}{EB9113}
\definecolor{quarto-callout-tip-color}{HTML}{00A047}
\definecolor{quarto-callout-caution-color}{HTML}{FC5300}
\definecolor{quarto-callout-color-frame}{HTML}{acacac}
\definecolor{quarto-callout-note-color-frame}{HTML}{4582ec}
\definecolor{quarto-callout-important-color-frame}{HTML}{d9534f}
\definecolor{quarto-callout-warning-color-frame}{HTML}{f0ad4e}
\definecolor{quarto-callout-tip-color-frame}{HTML}{02b875}
\definecolor{quarto-callout-caution-color-frame}{HTML}{fd7e14}
\makeatother
\makeatletter
\@ifpackageloaded{caption}{}{\usepackage{caption}}
\AtBeginDocument{%
\ifdefined\contentsname
  \renewcommand*\contentsname{Table of contents}
\else
  \newcommand\contentsname{Table of contents}
\fi
\ifdefined\listfigurename
  \renewcommand*\listfigurename{List of Figures}
\else
  \newcommand\listfigurename{List of Figures}
\fi
\ifdefined\listtablename
  \renewcommand*\listtablename{List of Tables}
\else
  \newcommand\listtablename{List of Tables}
\fi
\ifdefined\figurename
  \renewcommand*\figurename{Figure}
\else
  \newcommand\figurename{Figure}
\fi
\ifdefined\tablename
  \renewcommand*\tablename{Table}
\else
  \newcommand\tablename{Table}
\fi
}
\@ifpackageloaded{float}{}{\usepackage{float}}
\floatstyle{ruled}
\@ifundefined{c@chapter}{\newfloat{codelisting}{h}{lop}}{\newfloat{codelisting}{h}{lop}[chapter]}
\floatname{codelisting}{Listing}
\newcommand*\listoflistings{\listof{codelisting}{List of Listings}}
\makeatother
\makeatletter
\makeatother
\makeatletter
\@ifpackageloaded{caption}{}{\usepackage{caption}}
\@ifpackageloaded{subcaption}{}{\usepackage{subcaption}}
\makeatother

\usepackage{bookmark}
\IfFileExists{xurl.sty}{\usepackage{xurl}}{} % add URL line breaks if available
\urlstyle{same}
\hypersetup{
  pdftitle={Chapter 1: Loss Distributions},
  hidelinks,
  pdfcreator={LaTeX via pandoc}}


\title{\texorpdfstring{Chapter 1: Loss
Distributions}{Chapter 1: Loss Distributions}}
\author{}
\date{}

\begin{document}
\frame{\titlepage}


\section{Part I: Introduction}\label{part-i-introduction}

\begin{frame}{Loss Distributions}
\protect\phantomsection\label{loss-distributions}
\begin{block}{Learning Outcomes}
\protect\phantomsection\label{learning-outcomes}
By the end of this chapter, you should be able to

\begin{itemize}
\tightlist
\item
  explain why probability distributions are used to model insurance
  claim amounts;
\item
  distinguish between claim frequency and claim severity;
\item
  identify the common characteristics of insurance losses;
\item
  explain why different loss distributions are required for different
  types of insurance data; and
\item
  describe the situations in which the Exponential, Gamma, Lognormal and
  Pareto distributions are appropriate.
\end{itemize}
\end{block}
\end{frame}

\begin{frame}{Why Study Loss Distributions?}
\protect\phantomsection\label{why-study-loss-distributions}
Insurance transfers financial risk from policyholders to an insurer.

In exchange for a premium, the insurer agrees to compensate
policyholders when insured events occur.

Although the insurance contract is known in advance, \textbf{future
claim amounts are uncertain}.

The insurer therefore faces an important question:

\begin{quote}
\textbf{How much will future claims cost?}
\end{quote}

This question motivates the study of loss distributions.
\end{frame}

\begin{frame}{Why Is This Question Important?}
\protect\phantomsection\label{why-is-this-question-important}
Insurance companies must make financial decisions \textbf{before} future
claims are observed.

These decisions include

\begin{itemize}
\tightlist
\item
  determining premiums;
\item
  estimating claim liabilities;
\item
  purchasing reinsurance; and
\item
  holding sufficient capital to remain financially secure.
\end{itemize}

Because future losses are uncertain, insurers require mathematical
models to quantify this uncertainty.

\begin{tcolorbox}[enhanced jigsaw, arc=.35mm, bottomrule=.15mm, bottomtitle=1mm, breakable, colback=white, colbacktitle=quarto-callout-note-color!10!white, colframe=quarto-callout-note-color-frame, coltitle=black, left=2mm, leftrule=.75mm, opacityback=0, opacitybacktitle=0.6, rightrule=.15mm, title=\textcolor{quarto-callout-note-color}{\faInfo}\hspace{0.5em}{Key Idea}, titlerule=0mm, toprule=.15mm, toptitle=1mm]

Insurance decisions are made before future claims are observed.
Probability distributions provide a mathematical framework for
describing uncertainty in future claim amounts.

\end{tcolorbox}
\end{frame}

\begin{frame}{Motivating Example}
\protect\phantomsection\label{motivating-example}
Consider a motor insurer that issues \textbf{20,000 comprehensive motor
insurance policies}.

At the beginning of the year,

\begin{itemize}
\tightlist
\item
  the insurer knows the premium charged for each policy;
\item
  no one knows which policyholders will submit claims; and
\item
  no one knows how much those claims will cost.
\end{itemize}

By the end of the year,

\begin{itemize}
\tightlist
\item
  many policyholders make no claims;
\item
  many claims are relatively small;
\item
  a few claims are large; and
\item
  a very small number of claims are extremely large.
\end{itemize}

Although individual outcomes cannot be predicted, the overall pattern of
claims often remains remarkably stable from year to year.
\end{frame}

\begin{frame}{From Individual Claims to Statistical Models}
\protect\phantomsection\label{from-individual-claims-to-statistical-models}
Rather than predicting the outcome for a single policyholder, actuaries
study the behaviour of \textbf{large portfolios} of insurance policies.

Historical claims data reveal patterns that can be described using
probability distributions.

These models allow actuaries to estimate quantities such as

\begin{itemize}
\tightlist
\item
  the expected claim amount;
\item
  the variability of claim amounts; and
\item
  the probability of unusually large losses.
\end{itemize}

\begin{tcolorbox}[enhanced jigsaw, arc=.35mm, bottomrule=.15mm, bottomtitle=1mm, breakable, colback=white, colbacktitle=quarto-callout-note-color!10!white, colframe=quarto-callout-note-color-frame, coltitle=black, left=2mm, leftrule=.75mm, opacityback=0, opacitybacktitle=0.6, rightrule=.15mm, title=\textcolor{quarto-callout-note-color}{\faInfo}\hspace{0.5em}{Key Idea}, titlerule=0mm, toprule=.15mm, toptitle=1mm]

Actuarial models describe the overall behaviour of an insurance
portfolio rather than predicting individual claims.

\end{tcolorbox}
\end{frame}

\begin{frame}{Insurance Claims Data}
\protect\phantomsection\label{insurance-claims-data}
Every reported claim contributes to the insurer's claims database.

Typical information recorded for each claim includes

\begin{itemize}
\tightlist
\item
  policy type; date of loss; cause of the claim;
\item
  amount paid; and settlement status.
\end{itemize}

These data provide the foundation for statistical modelling.

\begin{columns}[T]
\begin{column}{0.48\linewidth}
{\def\LTcaptype{none} % do not increment counter
\begin{longtable}[]{@{}cr@{}}
\toprule\noalign{}
Policy & Claim Amount (USD) \\
\midrule\noalign{}
\endhead
1 & 0 \\
2 & 420 \\
3 & 1,350 \\
4 & 0 \\
5 & 18,750 \\
\bottomrule\noalign{}
\end{longtable}
}
\end{column}

\begin{column}{0.48\linewidth}
Several features are immediately apparent.

\begin{itemize}
\tightlist
\item
  Most policyholders make no claim.
\item
  Many claims are relatively small.
\item
  A few claims are much larger than the rest.
\end{itemize}

These features are typical of insurance claim data.
\end{column}
\end{columns}
\end{frame}

\begin{frame}{Claim Frequency and Claim Severity}
\protect\phantomsection\label{claim-frequency-and-claim-severity}
The insurer's total claims depend on \textbf{two distinct sources of
uncertainty}.

\begin{columns}[T]
\begin{column}{0.48\linewidth}
\begin{block}{Claim Frequency}
\protect\phantomsection\label{claim-frequency}
The \textbf{number of claims} occurring during a specified period.

Answers the question

\begin{quote}
\textbf{How many claims will occur?}
\end{quote}
\end{block}
\end{column}

\begin{column}{0.48\linewidth}
\begin{block}{Claim Severity}
\protect\phantomsection\label{claim-severity}
The \textbf{size of each individual claim}.

Answers the question

\begin{quote}
\textbf{How large will each claim be?}
\end{quote}
\end{block}
\end{column}
\end{columns}

Together, these determine the insurer's aggregate claims.
\end{frame}

\begin{frame}{Frequency and Severity}
\protect\phantomsection\label{frequency-and-severity}
\includegraphics[width=8.63in,height=1.81in]{01_Loss_Distributions_Slides_files/figure-beamer/mermaid-figure-1.png}

This chapter focuses on \textbf{claim severity}, that is, modelling the
probability distribution of individual claim amounts.

Models for claim frequency and aggregate claims are developed in later
chapters.
\end{frame}

\begin{frame}{Why Probability Distributions?}
\protect\phantomsection\label{why-probability-distributions}
Claim amounts vary from one policyholder to another.

Even when two policyholders experience similar accidents, the resulting
claim amounts may differ because of

\begin{itemize}
\tightlist
\item
  repair costs;
\item
  medical expenses;
\item
  policy conditions; and
\item
  inflation.
\end{itemize}

Since claim amounts are uncertain, they are naturally represented as
\textbf{random variables}.

A probability distribution describes

\begin{itemize}
\tightlist
\item
  the possible claim amounts; and
\item
  the probability associated with each amount.
\end{itemize}
\end{frame}

\begin{frame}{What Can a Probability Distribution Tell Us?}
\protect\phantomsection\label{what-can-a-probability-distribution-tell-us}
Once a suitable distribution has been selected, actuaries can calculate
quantities such as

\begin{itemize}
\tightlist
\item
  the expected claim amount;
\item
  the variability of losses;
\item
  the probability of large claims;
\item
  selected percentiles and risk measures; and
\item
  premiums based on expected losses.
\end{itemize}

Probability distributions therefore provide the mathematical foundation
for modelling insurance losses.
\end{frame}

\begin{frame}{Characteristics of Insurance Losses}
\protect\phantomsection\label{characteristics-of-insurance-losses}
Insurance claim amounts exhibit several common features.

\begin{itemize}
\tightlist
\item
  \textbf{Non-negative:} claim amounts cannot be negative.
\item
  \textbf{Right-skewed:} many small claims and relatively few large
  claims.
\item
  \textbf{Variable:} the spread of claim amounts differs across
  insurance classes.
\item
  \textbf{Heavy-tailed:} although rare, very large claims may have
  substantial financial consequences.
\end{itemize}

These characteristics strongly influence the choice of an appropriate
probability distribution.
\end{frame}

\begin{frame}{Characteristics of Insurance Losses}
\protect\phantomsection\label{characteristics-of-insurance-losses-1}
Insurance claim amounts exhibit several common features that influence
the choice of an appropriate loss distribution.

\begin{figure}[H]

{\centering \includegraphics[width=0.75\linewidth,height=\textheight,keepaspectratio]{../figures/Characteristics_of_Insurance_Losses.png}

}

\caption{Characteristics of Insurance Losses}

\end{figure}%
\end{frame}

\begin{frame}{Choosing an Appropriate Loss Distribution}
\protect\phantomsection\label{choosing-an-appropriate-loss-distribution}
No single probability distribution is suitable for every insurance
portfolio.

The choice of model depends on

\begin{itemize}
\tightlist
\item
  the characteristics of the claim data;
\item
  the degree of variability;
\item
  the behaviour of large claims; and
\item
  the purpose of the actuarial analysis.
\end{itemize}

Selecting an appropriate model is one of the first and most important
steps in actuarial modelling.
\end{frame}

\begin{frame}{Chapter Overview}
\protect\phantomsection\label{chapter-overview}
In this chapter we study four continuous distributions used to model
claim severity.

{\def\LTcaptype{none} % do not increment counter
\begin{longtable}[]{@{}ll@{}}
\toprule\noalign{}
Distribution & Typical Feature \\
\midrule\noalign{}
\endhead
Exponential & Simple model for non-negative losses \\
Gamma & More flexible than the Exponential distribution \\
Lognormal & Suitable for highly skewed claim amounts \\
Pareto & Models heavy-tailed insurance losses \\
\bottomrule\noalign{}
\end{longtable}
}

Together, these distributions provide increasingly flexible models for
insurance claim severity.
\end{frame}

\section{Part II: Exponential
Distribution}\label{part-ii-exponential-distribution}

\begin{frame}{Exponential Distribution}
\protect\phantomsection\label{exponential-distribution}
The Exponential distribution is the simplest continuous probability
distribution for modelling non-negative random variables.

Although many insurance claim amounts exhibit greater variability than
this model can capture, the Exponential distribution provides an
excellent starting point because many fundamental ideas in loss
modelling can be introduced using a single parameter.

Throughout this chapter, it serves as a foundation for understanding
more flexible distributions such as the Gamma, Lognormal and Pareto
distributions.

\begin{tcolorbox}[enhanced jigsaw, arc=.35mm, bottomrule=.15mm, bottomtitle=1mm, breakable, colback=white, colbacktitle=quarto-callout-note-color!10!white, colframe=quarto-callout-note-color-frame, coltitle=black, left=2mm, leftrule=.75mm, opacityback=0, opacitybacktitle=0.6, rightrule=.15mm, title=\textcolor{quarto-callout-note-color}{\faInfo}\hspace{0.5em}{Key Idea}, titlerule=0mm, toprule=.15mm, toptitle=1mm]

The Exponential distribution is a simple model for non-negative random
variables and introduces many of the concepts used throughout actuarial
loss modelling.

\end{tcolorbox}
\end{frame}

\begin{frame}{Why Study the Exponential Distribution?}
\protect\phantomsection\label{why-study-the-exponential-distribution}
The Exponential distribution is frequently used when

\begin{itemize}
\tightlist
\item
  observations are always non-negative;
\item
  small values occur much more frequently than large values; and
\item
  the probability decreases steadily as the variable increases.
\end{itemize}

Typical actuarial applications include

\begin{itemize}
\tightlist
\item
  time between successive insurance claims;
\item
  waiting time until the next event; and
\item
  relatively homogeneous claim amounts.
\end{itemize}

Although this model is often too simple for real claim severity data, it
provides the mathematical foundation for more advanced models.
\end{frame}

\begin{frame}{Definition of the Exponential Distribution}
\protect\phantomsection\label{definition-of-the-exponential-distribution}
A continuous random variable \(X\) is said to follow an
\textbf{Exponential distribution} with rate parameter \(\lambda>0\),
denoted by

\[
X\sim\operatorname{Exp}(\lambda),
\]

if its probability density function is

\[
f_X(x)=
\begin{cases}
\lambda e^{-\lambda x}, & x>0,\\
0, & \text{otherwise}.
\end{cases}
\]

The parameter \(\lambda\) is called the \textbf{rate parameter}.
\end{frame}

\begin{frame}{Understanding the Definition}
\protect\phantomsection\label{understanding-the-definition}
The probability density function satisfies the two requirements of a
valid density.

\begin{itemize}
\tightlist
\item
  It is always non-negative because \(\lambda>0\) and
  \(e^{-\lambda x}>0\).
\item
  The total probability equals one since
\end{itemize}

\[
\int_0^\infty \lambda e^{-\lambda x}\,dx=1.
\]

The support of the distribution is

\[
(0,\infty),
\]

which is appropriate for modelling quantities that cannot be negative,
such as insurance claim amounts.
\end{frame}

\begin{frame}{Interpreting the Rate Parameter}
\protect\phantomsection\label{interpreting-the-rate-parameter}
The rate parameter controls how rapidly the density decreases.

\begin{itemize}
\tightlist
\item
  A \textbf{large} value of \(\lambda\) causes the density to decrease
  rapidly.
\item
  Most observations are concentrated close to zero.
\item
  Large observations become relatively uncommon.
\end{itemize}

Conversely,

\begin{itemize}
\tightlist
\item
  a \textbf{small} value of \(\lambda\) produces a slower decline;
\item
  observations are spread over a wider range; and
\item
  larger values occur more frequently.
\end{itemize}
\end{frame}

\begin{frame}{The Rate Parameter and the Mean}
\protect\phantomsection\label{the-rate-parameter-and-the-mean}
The mean of an Exponential random variable is

\[
\mathrm{E}[X]=\frac1{\lambda}.
\]

This relationship provides a simple interpretation.

\begin{itemize}
\tightlist
\item
  Increasing \(\lambda\) decreases the average value.
\item
  Decreasing \(\lambda\) increases the average value.
\end{itemize}

\begin{tcolorbox}[enhanced jigsaw, arc=.35mm, bottomrule=.15mm, bottomtitle=1mm, breakable, colback=white, colbacktitle=quarto-callout-tip-color!10!white, colframe=quarto-callout-tip-color-frame, coltitle=black, left=2mm, leftrule=.75mm, opacityback=0, opacitybacktitle=0.6, rightrule=.15mm, title=\textcolor{quarto-callout-tip-color}{\faLightbulb}\hspace{0.5em}{Interpretation}, titlerule=0mm, toprule=.15mm, toptitle=1mm]

The rate parameter is inversely related to the average claim amount.

A larger rate implies a smaller average loss, while a smaller rate
implies a larger average loss.

\end{tcolorbox}
\end{frame}

\begin{frame}{Shape of the Density Function}
\protect\phantomsection\label{shape-of-the-density-function}
The Exponential density has several distinctive characteristics.

\begin{itemize}
\tightlist
\item
  The maximum density occurs at \(x=0\).
\item
  The density decreases continuously as \(x\) increases.
\item
  There is no interior mode.
\item
  The distribution is positively skewed.
\end{itemize}

This reflects situations where small claim amounts occur much more
frequently than large claim amounts.
\end{frame}

\begin{frame}{Tail Behaviour}
\protect\phantomsection\label{tail-behaviour}
Although the Exponential distribution has a right tail, the tail
decreases \textbf{exponentially}.

Consequently,

\begin{itemize}
\tightlist
\item
  moderately large observations may occur;
\item
  extremely large observations become increasingly unlikely.
\end{itemize}

This behaviour distinguishes the Exponential distribution from
heavy-tailed models such as the Lognormal and Pareto distributions,
which assign greater probability to very large claims.
\end{frame}

\begin{frame}{Density Function}
\protect\phantomsection\label{density-function}
The probability density function describes the relative likelihood of
different claim amounts.

Small claim amounts have the highest density, and the density decreases
continuously as the claim amount increases.

\begin{center}
\pandocbounded{\includegraphics[keepaspectratio]{01_Loss_Distributions_Slides_files/figure-beamer/unnamed-chunk-2-1.pdf}}
\end{center}
\end{frame}

\begin{frame}{Effect of the Rate Parameter}
\protect\phantomsection\label{effect-of-the-rate-parameter}
Changing \(\lambda\) changes the shape of the distribution.

\begin{itemize}
\tightlist
\item
  Larger \(\lambda\) values produce steeper curves.
\item
  Smaller \(\lambda\) values produce flatter curves with longer tails.
\end{itemize}

\begin{center}
\pandocbounded{\includegraphics[keepaspectratio]{01_Loss_Distributions_Slides_files/figure-beamer/unnamed-chunk-3-1.pdf}}
\end{center}
\end{frame}

\begin{frame}{Cumulative Distribution Function}
\protect\phantomsection\label{cumulative-distribution-function}
The cumulative distribution function gives the probability that a claim
amount does not exceed a specified value.

For an Exponential random variable,

\[
F_X(x)=
\begin{cases}
1-e^{-\lambda x}, & x>0,\\
0, & x\le0.
\end{cases}
\]

As \(x\) increases,

\begin{itemize}
\tightlist
\item
  the cumulative probability increases;
\item
  eventually it approaches one.
\end{itemize}
\end{frame}

\begin{frame}{Calculating Probabilities}
\protect\phantomsection\label{calculating-probabilities}
Once the cumulative distribution function is known, probabilities are
obtained directly.

For example,

\[
P(X\le x)=F_X(x)=1-e^{-\lambda x},
\]

while

\[
P(X>x)=1-F_X(x)=e^{-\lambda x}.
\]

The second expression is called the \textbf{survival function}, and it
plays an important role in insurance and reliability modelling.
\end{frame}

\begin{frame}{Mean and Variance}
\protect\phantomsection\label{mean-and-variance}
The first two moments are

\[
\mathrm{E}[X]=\frac1{\lambda},
\]

and

\[
\mathrm{Var}(X)=\frac1{\lambda^2}.
\]

These formulas show that both the average and the variability depend
entirely on the single parameter \(\lambda\).
\end{frame}

\begin{frame}{Moment Generating Function}
\protect\phantomsection\label{moment-generating-function}
The moment generating function is

\[
M_X(t)=
\frac{\lambda}{\lambda-t},
\qquad t<\lambda.
\]

The mgf is useful because

\begin{itemize}
\tightlist
\item
  moments can be obtained by differentiation; and
\item
  it provides an alternative way to study the distribution.
\end{itemize}
\end{frame}

\begin{frame}{Summary of the Exponential Distribution}
\protect\phantomsection\label{summary-of-the-exponential-distribution}
{\def\LTcaptype{none} % do not increment counter
\begin{longtable}[]{@{}ll@{}}
\toprule\noalign{}
Property & Result \\
\midrule\noalign{}
\endhead
Support & \((0,\infty)\) \\
Parameter & \(\lambda>0\) \\
Mean & \(\dfrac1{\lambda}\) \\
Variance & \(\dfrac1{\lambda^2}\) \\
PDF & \(\lambda e^{-\lambda x}\) \\
CDF & \(1-e^{-\lambda x}\) \\
MGF & \(\dfrac{\lambda}{\lambda-t},\;t<\lambda\) \\
\bottomrule\noalign{}
\end{longtable}
}
\end{frame}

\begin{frame}{Simulation from the Exponential Distribution}
\protect\phantomsection\label{simulation-from-the-exponential-distribution}
Random samples from an Exponential distribution can be generated easily
in R.
\end{frame}

\begin{frame}{Simulated Exponential Distribution}
\protect\phantomsection\label{simulated-exponential-distribution}
The figure below displays the simulated data overlayed with its true
theoretical density function.

\begin{center}
\pandocbounded{\includegraphics[keepaspectratio]{01_Loss_Distributions_Slides_files/figure-beamer/unnamed-chunk-4-1.pdf}}
\end{center}

The histogram approaches the theoretical density as the sample size
increases.
\end{frame}

\begin{frame}[fragile]{Simulated Exponential Distribution Code}
\protect\phantomsection\label{simulated-exponential-distribution-code}
\begin{columns}[T]
\begin{column}{0.48\linewidth}
\begin{block}{Part 1: Data Simulation}
\protect\phantomsection\label{part-1-data-simulation}
\begin{Shaded}
\begin{Highlighting}[]
\CommentTok{\# 1. Load required libraries}
\FunctionTok{library}\NormalTok{(ggplot2)}

\CommentTok{\# 2. Set seed for reproducibility}
\FunctionTok{set.seed}\NormalTok{(}\DecValTok{123}\NormalTok{)}

\CommentTok{\# 3. Simulate exponential data}
\NormalTok{df }\OtherTok{\textless{}{-}} \FunctionTok{data.frame}\NormalTok{(}
  \AttributeTok{x =} \FunctionTok{rexp}\NormalTok{(}\DecValTok{1000}\NormalTok{, }\AttributeTok{rate =} \FloatTok{0.2}\NormalTok{)}
\NormalTok{)}
\end{Highlighting}
\end{Shaded}
\end{block}
\end{column}

\begin{column}{0.04\linewidth}
\end{column}

\begin{column}{0.48\linewidth}
\begin{block}{Part 2: Density Plotting}
\protect\phantomsection\label{part-2-density-plotting}
\begin{Shaded}
\begin{Highlighting}[]
\CommentTok{\# 4. Initialize ggplot setup}
\FunctionTok{ggplot}\NormalTok{(df, }\FunctionTok{aes}\NormalTok{(}\AttributeTok{x =}\NormalTok{ x)) }\SpecialCharTok{+}
  \FunctionTok{geom\_histogram}\NormalTok{(}
    \FunctionTok{aes}\NormalTok{(}\AttributeTok{y =} \FunctionTok{after\_stat}\NormalTok{(density)), }
    \AttributeTok{bins =} \DecValTok{30}\NormalTok{, }
    \AttributeTok{color =} \StringTok{"white"}\NormalTok{, }
    \AttributeTok{fill =} \StringTok{"steelblue"}
\NormalTok{  ) }\SpecialCharTok{+}
  \CommentTok{\# 5. Overlay theoretical density curve}
  \FunctionTok{stat\_function}\NormalTok{(}
    \AttributeTok{fun =}\NormalTok{ dexp, }
    \AttributeTok{args =} \FunctionTok{list}\NormalTok{(}\AttributeTok{rate =} \FloatTok{0.2}\NormalTok{), }
    \AttributeTok{color =} \StringTok{"black"}\NormalTok{, }
    \AttributeTok{linewidth =} \FloatTok{0.8}
\NormalTok{  ) }\SpecialCharTok{+}
  \FunctionTok{labs}\NormalTok{(}
    \AttributeTok{y =} \StringTok{"Density"}\NormalTok{, }
    \AttributeTok{x =} \StringTok{"Claim amount"}
\NormalTok{  ) }\SpecialCharTok{+}
  \FunctionTok{theme\_minimal}\NormalTok{()}
\end{Highlighting}
\end{Shaded}
\end{block}
\end{column}
\end{columns}
\end{frame}

\begin{frame}{Looking Ahead}
\protect\phantomsection\label{looking-ahead}
The Exponential distribution is mathematically simple and provides a
useful introduction to continuous loss models.

However, many insurance claim datasets exhibit greater variability than
this distribution can accommodate.

The next distribution, the \textbf{Gamma distribution}, extends the
Exponential model by introducing an additional parameter, providing
substantially greater flexibility while retaining many of the same
mathematical properties.
\end{frame}

\section{Part III: Gamma
Distribution}\label{part-iii-gamma-distribution}

\begin{frame}{Gamma Distribution}
\protect\phantomsection\label{gamma-distribution}
The Gamma distribution extends the Exponential distribution by
introducing a second parameter.

This additional parameter allows the distribution to model a much wider
range of claim amount data while retaining many of the mathematical
properties of the Exponential distribution.

\begin{tcolorbox}[enhanced jigsaw, arc=.35mm, bottomrule=.15mm, bottomtitle=1mm, breakable, colback=white, colbacktitle=quarto-callout-note-color!10!white, colframe=quarto-callout-note-color-frame, coltitle=black, left=2mm, leftrule=.75mm, opacityback=0, opacitybacktitle=0.6, rightrule=.15mm, title=\textcolor{quarto-callout-note-color}{\faInfo}\hspace{0.5em}{Key Idea}, titlerule=0mm, toprule=.15mm, toptitle=1mm]

The Exponential distribution is a special case of the Gamma
distribution.

\end{tcolorbox}
\end{frame}

\begin{frame}{Why Introduce the Gamma Distribution?}
\protect\phantomsection\label{why-introduce-the-gamma-distribution}
The Exponential distribution assumes that the mean and variability are
both determined by a single parameter.

In practice, insurance claim amounts often exhibit greater variability
than this simple model can capture.

The Gamma distribution introduces an additional \textbf{shape
parameter}, allowing greater flexibility while remaining mathematically
tractable.
\end{frame}

\begin{frame}{Definition}
\protect\phantomsection\label{definition}
A random variable \(X\) has a Gamma distribution with shape parameter
\(\alpha>0\) and rate parameter \(\lambda>0\), denoted by

\[
X\sim\operatorname{Gamma}(\alpha,\lambda),
\]

if its probability density function is

\[
f_X(x)=
\frac{\lambda^\alpha}{\Gamma(\alpha)}
x^{\alpha-1}
e^{-\lambda x},
\qquad x>0.
\]
\end{frame}

\begin{frame}{The Gamma Function}
\protect\phantomsection\label{the-gamma-function}
The Gamma function is defined by

\[
\Gamma(\alpha)
=
\int_0^\infty
x^{\alpha-1}e^{-x}\,dx,
\qquad
\alpha>0.
\]

It generalises the factorial function since

\[
\Gamma(n)=(n-1)!,
\qquad n=1,2,\ldots
\]

This ensures that the Gamma density integrates to one.
\end{frame}

\begin{frame}{Relationship with the Exponential Distribution}
\protect\phantomsection\label{relationship-with-the-exponential-distribution}
When

\[
\alpha=1,
\]

the Gamma density becomes

\[
f(x)
=
\lambda e^{-\lambda x},
\]

which is exactly the Exponential distribution.

Thus,

\[
\operatorname{Exp}(\lambda)
=
\Gamma(1,\lambda).
\]
\end{frame}

\begin{frame}{Distribution Function}
\protect\phantomsection\label{distribution-function}
The cumulative distribution function is

\[
F_X(x)
=
\int_0^x
f_X(t)\,dt.
\]

Unlike the Exponential distribution, no simple closed-form expression
exists for general values of \(\alpha\).

In practice, probabilities are evaluated using statistical software.
\end{frame}

\begin{frame}{Properties}
\protect\phantomsection\label{properties}
{\def\LTcaptype{none} % do not increment counter
\begin{longtable}[]{@{}
  >{\raggedright\arraybackslash}p{(\linewidth - 2\tabcolsep) * \real{0.5789}}
  >{\raggedright\arraybackslash}p{(\linewidth - 2\tabcolsep) * \real{0.4211}}@{}}
\toprule\noalign{}
\begin{minipage}[b]{\linewidth}\raggedright
Property
\end{minipage} & \begin{minipage}[b]{\linewidth}\raggedright
Result
\end{minipage} \\
\midrule\noalign{}
\endhead
Mean & \(\dfrac{\alpha}{\lambda}\) \\
Variance & \(\dfrac{\alpha}{\lambda^2}\) \\
MGF &
\(\left(\dfrac{\lambda}{\lambda-t}\right)^\alpha,\quad t<\lambda\) \\
\bottomrule\noalign{}
\end{longtable}
}

Because two parameters are available, the mean and variance can be
adjusted more flexibly than for the Exponential distribution.
\end{frame}

\begin{frame}{Effect of the Shape Parameter}
\protect\phantomsection\label{effect-of-the-shape-parameter}
The shape parameter controls the overall form of the distribution.

\begin{itemize}
\tightlist
\item
  \(\alpha=1\) gives the Exponential distribution.
\item
  Larger values of \(\alpha\) produce a more symmetric density.
\item
  Small values of \(\alpha\) produce a strongly right-skewed
  distribution.
\end{itemize}
\end{frame}

\begin{frame}{Shape of the Density}
\protect\phantomsection\label{shape-of-the-density}
The figure below illustrates how the density changes for different
parameter values.

\begin{center}
\pandocbounded{\includegraphics[keepaspectratio]{01_Loss_Distributions_Slides_files/figure-beamer/unnamed-chunk-7-1.pdf}}
\end{center}
\end{frame}

\begin{frame}{When Is the Gamma Distribution Useful?}
\protect\phantomsection\label{when-is-the-gamma-distribution-useful}
The Gamma distribution is widely used for modelling

\begin{itemize}
\tightlist
\item
  insurance claim amounts,
\item
  repair costs,
\item
  waiting times,
\item
  aggregate losses.
\end{itemize}

Compared with the Exponential distribution, it provides a better fit
when claim data exhibit moderate variability.
\end{frame}

\begin{frame}{Summary}
\protect\phantomsection\label{summary}
Compared with the Exponential distribution, the Gamma distribution

\begin{itemize}
\tightlist
\item
  introduces an additional shape parameter;
\item
  provides substantially greater modelling flexibility;
\item
  includes the Exponential distribution as a special case; and
\item
  remains one of the most widely used continuous loss distributions in
  actuarial science.
\end{itemize}
\end{frame}

\begin{frame}{Looking Ahead}
\protect\phantomsection\label{looking-ahead-1}
Although the Gamma distribution is highly flexible, some insurance claim
data exhibit even stronger positive skewness and heavier tails.

The next distribution, the \textbf{Lognormal distribution}, is often
more appropriate for modelling these types of losses.
\end{frame}

\section{Part IV: Lognormal
Distribution}\label{part-iv-lognormal-distribution}

\begin{frame}{Lognormal Distribution}
\protect\phantomsection\label{lognormal-distribution}
Many insurance claim amounts are highly right-skewed.

Although the Gamma distribution provides greater flexibility than the
Exponential distribution, it may still underestimate the probability of
very large claims.

The Lognormal distribution is often a better choice for modelling claim
amounts with a long right tail.

\begin{tcolorbox}[enhanced jigsaw, arc=.35mm, bottomrule=.15mm, bottomtitle=1mm, breakable, colback=white, colbacktitle=quarto-callout-note-color!10!white, colframe=quarto-callout-note-color-frame, coltitle=black, left=2mm, leftrule=.75mm, opacityback=0, opacitybacktitle=0.6, rightrule=.15mm, title=\textcolor{quarto-callout-note-color}{\faInfo}\hspace{0.5em}{Key Idea}, titlerule=0mm, toprule=.15mm, toptitle=1mm]

A random variable is \textbf{Lognormally distributed} if its logarithm
follows a Normal distribution.

\end{tcolorbox}
\end{frame}

\begin{frame}{Why Introduce the Lognormal Distribution?}
\protect\phantomsection\label{why-introduce-the-lognormal-distribution}
The Lognormal distribution is appropriate when

\begin{itemize}
\tightlist
\item
  claim amounts are always positive;
\item
  most claims are relatively small;
\item
  a small number of claims are extremely large.
\end{itemize}

These characteristics are commonly observed in insurance claim severity
data.
\end{frame}

\begin{frame}{Definition}
\protect\phantomsection\label{definition-1}
A random variable \(X\) has a \textbf{Lognormal distribution} with
parameters \(\mu\) and \(\sigma^2\), denoted by

\[
X \sim \operatorname{LN}(\mu,\sigma^2),
\]

if

\[
\ln(X)\sim N(\mu,\sigma^2).
\]

Thus, the logarithm of the claim amount follows a Normal distribution.
\end{frame}

\begin{frame}{Probability Density Function}
\protect\phantomsection\label{probability-density-function}
The probability density function is

\[
f_X(x)=
\frac{1}{x\sigma\sqrt{2\pi}}
\exp\!\left[
-\frac{(\ln x-\mu)^2}{2\sigma^2}
\right],
\qquad x>0.
\]

The support is

\[
(0,\infty),
\]

making the Lognormal distribution suitable for modelling positive claim
amounts.
\end{frame}

\begin{frame}{Distribution Function}
\protect\phantomsection\label{distribution-function-1}
The cumulative distribution function is

\[
F_X(x)
=
\Phi
\left(
\frac{\ln x-\mu}{\sigma}
\right),
\qquad x>0,
\]

where \(\Phi(\cdot)\) denotes the cumulative distribution function of
the standard Normal distribution.

In practice, probabilities are evaluated using statistical software.
\end{frame}

\begin{frame}{Properties}
\protect\phantomsection\label{properties-1}
{\def\LTcaptype{none} % do not increment counter
\begin{longtable}[]{@{}ll@{}}
\toprule\noalign{}
Property & Result \\
\midrule\noalign{}
\endhead
Mean & \(e^{\mu+\sigma^2/2}\) \\
Variance & \(e^{2\mu+\sigma^2}\left(e^{\sigma^2}-1\right)\) \\
MGF & Does not exist for \(t>0\) \\
\bottomrule\noalign{}
\end{longtable}
}

Unlike the Exponential and Gamma distributions, the Lognormal
distribution has no finite moment generating function.
\end{frame}

\begin{frame}{Understanding the Parameters}
\protect\phantomsection\label{understanding-the-parameters}
The two parameters influence the distribution in different ways.

\begin{itemize}
\tightlist
\item
  \(\mu\) controls the location of the distribution.
\item
  \(\sigma\) controls the spread and the degree of right skewness.
\end{itemize}

As \(\sigma\) increases,

\begin{itemize}
\tightlist
\item
  the distribution becomes more dispersed;
\item
  the right tail becomes longer; and
\item
  extremely large claims become more likely.
\end{itemize}
\end{frame}

\begin{frame}{Shape of the Density}
\protect\phantomsection\label{shape-of-the-density-1}
Compared with the Gamma distribution, the Lognormal density generally
has

\begin{itemize}
\tightlist
\item
  stronger positive skewness;
\item
  a longer right tail; and
\item
  greater probability of very large losses.
\end{itemize}
\end{frame}

\begin{frame}{Shape of the Density}
\protect\phantomsection\label{shape-of-the-density-2}
The figure below illustrates how the density changes for different
parameter values.

\begin{center}
\pandocbounded{\includegraphics[keepaspectratio]{01_Loss_Distributions_Slides_files/figure-beamer/unnamed-chunk-8-1.pdf}}
\end{center}
\end{frame}

\begin{frame}{When Is the Lognormal Distribution Useful?}
\protect\phantomsection\label{when-is-the-lognormal-distribution-useful}
The Lognormal distribution is widely used to model

\begin{itemize}
\tightlist
\item
  motor insurance claim amounts;
\item
  property insurance losses;
\item
  medical expenses;
\item
  liability claims.
\end{itemize}

It often provides a better fit than the Exponential or Gamma
distributions when large claims occur more frequently.
\end{frame}

\begin{frame}{Summary}
\protect\phantomsection\label{summary-1}
The Lognormal distribution

\begin{itemize}
\tightlist
\item
  models positive and highly skewed claim amounts;
\item
  is defined through the Normal distribution;
\item
  produces a heavier right tail than the Exponential and Gamma
  distributions; and
\item
  is widely used for modelling claim severity.
\end{itemize}
\end{frame}

\begin{frame}{Looking Ahead}
\protect\phantomsection\label{looking-ahead-2}
Although the Lognormal distribution models heavy-tailed losses well,
some insurance portfolios contain extremely large claims that require an
even heavier-tailed model.

The next section introduces the \textbf{Pareto distribution}, which is
widely used in reinsurance and extreme loss modelling.
\end{frame}

\section{Part V: Pareto Distribution}\label{part-v-pareto-distribution}

\begin{frame}{Pareto Distribution}
\protect\phantomsection\label{pareto-distribution}
The Exponential, Gamma and Lognormal distributions are useful for
modelling many insurance claim amounts.

However, some portfolios contain \textbf{very large claims} that occur
more frequently than these models predict.

The Pareto distribution provides a simple model for such
\textbf{heavy-tailed} losses.

\begin{tcolorbox}[enhanced jigsaw, arc=.35mm, bottomrule=.15mm, bottomtitle=1mm, breakable, colback=white, colbacktitle=quarto-callout-note-color!10!white, colframe=quarto-callout-note-color-frame, coltitle=black, left=2mm, leftrule=.75mm, opacityback=0, opacitybacktitle=0.6, rightrule=.15mm, title=\textcolor{quarto-callout-note-color}{\faInfo}\hspace{0.5em}{Key Idea}, titlerule=0mm, toprule=.15mm, toptitle=1mm]

The Pareto distribution assigns relatively more probability to extremely
large losses than the Exponential, Gamma and Lognormal distributions.

\end{tcolorbox}
\end{frame}

\begin{frame}{Why Introduce the Pareto Distribution?}
\protect\phantomsection\label{why-introduce-the-pareto-distribution}
The Pareto distribution is widely used for modelling

\begin{itemize}
\tightlist
\item
  liability claims;
\item
  catastrophe losses;
\item
  excess-of-loss reinsurance claims;
\item
  large commercial insurance losses.
\end{itemize}

These applications involve claim amounts with long right tails.
\end{frame}

\begin{frame}{Definition}
\protect\phantomsection\label{definition-2}
A random variable \(X\) has a Pareto distribution with parameters
\(\alpha>0\) and \(\lambda>0\), denoted by

\[
X\sim\operatorname{Pa}(\alpha,\lambda),
\]

if its probability density function is

\[
f_X(x)
=
\frac{\alpha\lambda^\alpha}
{(\lambda+x)^{\alpha+1}},
\qquad x>0.
\]

The cumulative distribution function is

\[
F_X(x)
=
1-
\left(
\frac{\lambda}{\lambda+x}
\right)^\alpha,
\qquad x>0.
\]
\end{frame}

\begin{frame}{Properties}
\protect\phantomsection\label{properties-2}
{\def\LTcaptype{none} % do not increment counter
\begin{longtable}[]{@{}
  >{\raggedright\arraybackslash}p{(\linewidth - 2\tabcolsep) * \real{0.5789}}
  >{\raggedright\arraybackslash}p{(\linewidth - 2\tabcolsep) * \real{0.4211}}@{}}
\toprule\noalign{}
\begin{minipage}[b]{\linewidth}\raggedright
Property
\end{minipage} & \begin{minipage}[b]{\linewidth}\raggedright
Result
\end{minipage} \\
\midrule\noalign{}
\endhead
Mean & \(\displaystyle \frac{\lambda}{\alpha-1},\;\alpha>1\) \\
Variance &
\(\displaystyle \frac{\alpha\lambda^2}{(\alpha-1)^2(\alpha-2)},\;\alpha>2\) \\
\(r\)th moment &
\(\displaystyle \frac{\Gamma(\alpha-r)\Gamma(r+1)}{\Gamma(\alpha)}\lambda^r,\;0<r<\alpha\) \\
MGF & Does not exist for \(t>0\) \\
\bottomrule\noalign{}
\end{longtable}
}

The existence of the moments depends on the value of the shape parameter
\(\alpha\).
\end{frame}

\begin{frame}{Understanding the Parameters}
\protect\phantomsection\label{understanding-the-parameters-1}
The two parameters play different roles.

\begin{itemize}
\tightlist
\item
  \(\alpha\) controls the heaviness of the tail.
\item
  \(\lambda\) controls the scale of the distribution.
\end{itemize}

Smaller values of \(\alpha\) produce heavier tails and increase the
probability of extremely large claims.
\end{frame}

\begin{frame}{Shape of the Density}
\protect\phantomsection\label{shape-of-the-density-3}
The figure below illustrates how the density changes for different
parameter values.

\begin{center}
\pandocbounded{\includegraphics[keepaspectratio]{01_Loss_Distributions_Slides_files/figure-beamer/unnamed-chunk-9-1.pdf}}
\end{center}
\end{frame}

\begin{frame}{Conditional Tail Property}
\protect\phantomsection\label{conditional-tail-property}
One of the most useful properties of the Pareto distribution is its
conditional tail behaviour.

If

\[
X\sim\operatorname{Pa}(\alpha,\lambda),
\]

then

\[
X-w\mid X>w
\sim
\operatorname{Pa}(\alpha,\lambda+w).
\]

This property greatly simplifies calculations involving excess losses.
\end{frame}

\begin{frame}{Application to Reinsurance}
\protect\phantomsection\label{application-to-reinsurance}
In excess-of-loss reinsurance,

the insurer is responsible only for the amount exceeding a retention
level \(w\).

Because the conditional excess remains Pareto distributed,

the expected payment above the retention can often be obtained
analytically.

This is one reason why the Pareto distribution is widely used in
actuarial modelling.
\end{frame}

\begin{frame}{Comparing the Four Distributions}
\protect\phantomsection\label{comparing-the-four-distributions}
{\def\LTcaptype{none} % do not increment counter
\begin{longtable}[]{@{}lll@{}}
\toprule\noalign{}
Distribution & Tail Behaviour & Typical Use \\
\midrule\noalign{}
\endhead
Exponential & Light & Simple claim models \\
Gamma & Moderate & General claim severity \\
Lognormal & Heavy & Skewed claim amounts \\
Pareto & Very heavy & Large and catastrophic claims \\
\bottomrule\noalign{}
\end{longtable}
}

The Lognormal and Pareto distributions generally provide a better fit to
claim severity data than the Exponential and Gamma distributions.
\end{frame}

\begin{frame}{Other Heavy-Tailed Models}
\protect\phantomsection\label{other-heavy-tailed-models}
The Pareto distribution is one member of a broader family of loss
models.

Other distributions commonly used in actuarial science include

\begin{itemize}
\tightlist
\item
  Burr distribution,
\item
  Weibull distribution,
\item
  Loggamma distribution.
\end{itemize}

These models provide additional flexibility when fitting insurance claim
data.
\end{frame}

\begin{frame}{Example: Fitting Loss Distributions}
\protect\phantomsection\label{example-fitting-loss-distributions}
Consider a portfolio containing \textbf{200 observed claim amounts}.

The objective is to compare several candidate loss distributions.

The analysis consists of the following steps.

\begin{enumerate}
\tightlist
\item
  Calculate the sample mean and sample standard deviation.
\item
  Fit Exponential and Gamma distributions using the method of moments.
\item
  Construct class intervals with expected frequency 20 under the fitted
  Exponential distribution.
\item
  Count the observed claims in each interval.
\item
  Compute the expected frequencies under the Gamma, Lognormal and Pareto
  models.
\item
  Compare histograms with the fitted density curves.
\end{enumerate}
\end{frame}

\begin{frame}{Discussion}
\protect\phantomsection\label{discussion}
The fitted models can be compared by examining

\begin{itemize}
\tightlist
\item
  agreement between observed and expected frequencies;
\item
  goodness of fit of the density curves;
\item
  ability to represent large claims.
\end{itemize}

In practice,

the Lognormal and Pareto distributions often provide a substantially
better fit than the Exponential and Gamma distributions because they
assign greater probability to large claim amounts.
\end{frame}

\begin{frame}{Chapter Summary}
\protect\phantomsection\label{chapter-summary}
This chapter introduced four continuous loss distributions.

\begin{itemize}
\tightlist
\item
  Exponential: simplest one-parameter model.
\item
  Gamma: more flexible two-parameter model.
\item
  Lognormal: suitable for highly skewed losses.
\item
  Pareto: suitable for heavy-tailed and extreme claims.
\end{itemize}

These distributions provide the foundation for modelling claim severity
in actuarial science and prepare us for the study of deductibles,
reinsurance and collective risk models.
\end{frame}




\end{document}
